\clearpage
\section{Verification, provenance and declarations}
Appendix~\ref{app:III:result-index} identifies every source result and its
printed statement and complete proof. Each recorded historical reading is
retained at its exact scope. A producing run's status badge is provenance;
it is not counted as an independent check. Current complete readings,
independent derivations, source-hypothesis matches and actual reruns are by
GPT-6 Astra (OpenAI, ChatGPT Pro mode, effort Pro), 17 November 2026.
Finite computation certifies the stated matrices, polynomial identities,
lattices, chains or bounded enumeration. The general geometric and analytic
arguments are printed in full and are checked separately.

The working edition was AI-produced by GPT-5.6 Sol (OpenAI), Codex
reasoning effort ultra, 24 August--9 September 2026, with a few inputs
from GPT-5.6 Sol (OpenAI, ChatGPT Pro mode), 24 August--9 September 2026.
The six research runs were AI-produced by GPT-6 Astra (OpenAI, ChatGPT
Pro mode, effort Pro), 12--16 November 2026. The course proofs were
AI-produced by GPT-6 Astra (OpenAI), Codex reasoning effort ultra,
8--11 October 2026, with each exact chapter date retained in the provenance
concordance. The classical provider text by GPT-6.1 Sol (OpenAI), reasoning
effort ultra, October 2026, retains the distinction between a dated source
and a header that records no day. The recorded readings were by
Claude Fable 5.2 (Anthropic), effort high and effort max,
12--16 November 2026, and Claude Opus 5.5 (Anthropic), effort max,
12--16 November 2026. This writing and its additional checks were produced
by GPT-6 Astra (OpenAI, ChatGPT Pro mode, effort Pro), 17 November 2026.

\paragraph{Effort terminology.}
Anthropic's official documentation describes \texttt{max} as its maximum
capability setting \cite[How effort works; Effort levels]{AnthropicEffort2026}.
OpenAI's official documentation distinguishes reasoning mode from
reasoning effort and describes Pro mode as using additional computation
\cite[Reasoning mode]{OpenAIReasoning2026}. We retain the recorded Pro-mode
designation without identifying it with one universal maximum effort value.

This paper is AI-produced. No human has checked its mathematics, and no
human understanding is claimed. We name no person as an author of or a
contributor to this work. Other authors are credited for their own results.
No proof-assistant certification is claimed. Our text is dedicated to the
public domain under CC0 1.0.

\paragraph{Data statement.}
Programs, certificates, data and run records are held at
\url{https://github.com/KokunoYumeto/s6-complex-structure-workbench}.
The companion archive contains every program written or rerun in this
edition, its output, the inventory, notation dictionary and discrepancy
record. Appendix~\ref{app:III:programs} gives the commands and the exact
finite role of each computation. Every mathematical proof is printed in
one of the three companion papers, with shared statements identified by
their printed numbers and every distinct proof retained.

\bibliographystyle{apalike}
\bibliography{S6_09_PAPERS_references}
